\documentclass[11pt,a4paper]{article}
\usepackage[margin=27mm,headheight=15pt]{geometry}
\usepackage[T1]{fontenc}
\usepackage[utf8]{inputenc}
\usepackage{amsmath,amsthm}
\usepackage{newpxtext,newpxmath}
\usepackage{microtype}
\usepackage{mathtools}
\usepackage{booktabs,array,longtable}
\usepackage{xcolor}
\usepackage{enumitem,aliascnt}
\usepackage{fancyhdr}
\usepackage{titlesec}
\usepackage{listings}
\usepackage[most]{tcolorbox}
\definecolor{OliveInk}{HTML}{46523B}
\definecolor{MutedInk}{HTML}{62675D}
\definecolor{SoftOlive}{HTML}{F1F3EB}
\definecolor{RuleOlive}{HTML}{B6BEA5}
\usepackage[colorlinks=true,linkcolor=OliveInk,citecolor=OliveInk,urlcolor=OliveInk,
 pdfauthor={Research report prepared for Vladimir Reshetnikov},
 pdftitle={A sharp order-five threshold for log-concavity of higher-order Stirling subset numbers}]{hyperref}
\usepackage[nameinlink,noabbrev]{cleveref}
\titleformat{\section}{\normalfont\Large\bfseries\color{OliveInk}}{\thesection}{0.6em}{}
\titleformat{\subsection}{\normalfont\large\bfseries\color{OliveInk}}{\thesubsection}{0.6em}{}
\pagestyle{fancy}
\fancyhf{}
\fancyhead[L]{\small\color{MutedInk}Higher-order Stirling log-concavity}
\fancyhead[R]{\small\color{MutedInk}Research manuscript}
\fancyfoot[C]{\small\thepage}
\renewcommand{\headrulewidth}{0.3pt}
\setlength{\parskip}{0.32em}
\setlength{\parindent}{1.25em}
\setlength{\emergencystretch}{2em}
\setcounter{tocdepth}{1}
\numberwithin{equation}{section}
\newtheorem{theorem}{Theorem}[section]
\newaliascnt{lemma}{theorem}
\newtheorem{lemma}[lemma]{Lemma}
\aliascntresetthe{lemma}
\crefname{lemma}{lemma}{lemmas}
\newaliascnt{proposition}{theorem}
\newtheorem{proposition}[proposition]{Proposition}
\aliascntresetthe{proposition}
\crefname{proposition}{proposition}{propositions}
\newaliascnt{corollary}{theorem}
\newtheorem{corollary}[corollary]{Corollary}
\aliascntresetthe{corollary}
\crefname{corollary}{corollary}{corollaries}
\theoremstyle{definition}
\newaliascnt{definition}{theorem}
\newtheorem{definition}[definition]{Definition}
\aliascntresetthe{definition}
\crefname{definition}{definition}{definitions}
\newaliascnt{example}{theorem}
\newtheorem{example}[example]{Example}
\aliascntresetthe{example}
\crefname{example}{example}{examples}
\theoremstyle{remark}
\newaliascnt{remark}{theorem}
\newtheorem{remark}[remark]{Remark}
\aliascntresetthe{remark}
\crefname{remark}{remark}{remarks}
\newcommand{\Z}{\mathbb Z}
\newcommand{\R}{\mathbb R}
\newcommand{\Q}{\mathbb Q}
\newcommand{\N}{\mathbb N}
\newcommand{\lc}{\operatorname{LC}}
\newcommand{\Snum}[3]{S^{(#1)}_{#2,#3}}
\newcommand{\Cnum}[3]{C^{(#1)}_{#2,#3}}
\newcommand{\assocS}[3]{\genfrac\{\}{0pt}{}{#1}{#2}_{\!\ge #3}}
\newcommand{\assocC}[3]{\genfrac[]{0pt}{}{#1}{#2}_{\!\ge #3}}
\newcommand{\Turan}{Tur\'an}
\newcommand{\dup}[1]{{\small\textsf{\color{MutedInk}[#1]}}}
\lstset{basicstyle=\small\ttfamily,columns=fullflexible,keepspaces=true,
 breaklines=true,frame=single,rulecolor=\color{RuleOlive},
 backgroundcolor=\color{SoftOlive},showstringspaces=false,tabsize=4}
\tcbset{colback=SoftOlive,colframe=RuleOlive,boxrule=0.5pt,arc=1mm,
 left=3mm,right=3mm,top=2mm,bottom=2mm}

\begin{document}
\begin{center}
{\small\scshape\color{MutedInk}Combinatorics \quad / \quad Integer sequences \quad / \quad Recurrences}\par
\vspace{0.6cm}
{\LARGE\bfseries\color{OliveInk}A sharp order-five threshold for\par
log-concavity of higher-order\par Stirling subset numbers}\par
\vspace{0.45cm}
{\large An elementary recurrence proof, an exact fifth-order certificate,\par
and applications to Stirling cycle triangles}\par
\vspace{0.45cm}
{\small Merged research report prepared for Vladimir Reshetnikov\par
19--20 September 2026}
\end{center}
\vspace{0.3cm}
\begin{abstract}
For an integer $r\ge1$, let $S^{(r)}_{n,k}$ count partitions of an
$(n+(r-1)k)$-element set into $k$ blocks, each of size at least $r$.
We prove that all rows of this diagonal-shifted Stirling triangle are
log-concave if and only if $1\le r\le5$, and that for those orders every
internal inequality with $n\ge2$ and $1\le k\le n-1$ is strict. This
establishes the log-concavity assertion of Deb--Sokal's
Conjecture~1.4(c), including its cases $r=3,4,5$, and nothing else of
that conjecture. The proof reduces a two-term recurrence to a quadratic
form. For $r=5$ the decisive step is the exact identity
$25(A_4u^2+H_4uv+v^2)=(5v-8\alpha u)^2+96(t-5)\beta\gamma u^2
+480\beta(t-5k)uv$, whose three summands are nonnegative precisely on the
triangular domain; it also supplies an explicit lower bound for each
\Turan{} difference. The converse follows from an exact counterexample at
$(r,n,k)=(6,4,2)$, derived here in two independent ways, and a monotone
factorial ratio that gives counterexamples in row three for every $r\ge7$.
The same method proves log-concavity for the cycle triangles of orders
three and four and for a continuous family of interpolating recurrences,
and an asymptotic analysis explains why the elementary mechanism changes
exactly between orders five and six. We also exhibit a precise obstruction
to extending the method to the fifth-order cycle triangle; that separate
conjectural case is not settled here. All arguments are given in ordinary
mathematical form. Two independent verification toolchains accompany the
report: a standard-library package that checks the polynomial identities
exactly without a computer algebra system, and a SymPy checker of the same
identities.
\end{abstract}

\begin{tcolorbox}
\textbf{Scope and verification.} This is an unrefereed research manuscript,
not an independently certified publication. The selected subset-number
assertion has a complete proof below; the claim is not based on a numerical
search. The intended resolution is precisely the \emph{log-concavity clause}
of Conjecture~1.4(c) in \cite{DebSokal}; it is not a resolution of that
paper's total-positivity, real-rootedness, or cycle-number conjectures.
The originating preprint still states the selected assertion as a conjecture
in the version accessed for this report. A targeted literature check did not
locate a later resolution, but this does not certify priority. This
manuscript has not undergone independent peer review or proof-assistant
verification.
\end{tcolorbox}
\newpage
\tableofcontents
\vspace{0.9cm}
\begin{tcolorbox}[title={\bfseries Reading guide},colbacktitle=SoftOlive,
 coltitle=OliveInk]
The proof of the selected conjecture is in Sections~2--6.
Section~3 isolates the local inequality; Section~4 contains the
order-five certificate. A reviewer in a hurry should read the single
identity~\eqref{eq:SOS}, then the six-term display~\eqref{eq:full-certificate},
which is the entire interior induction step in one equation.
Sections~7--9 give the cycle-number extensions, an exact obstruction, and
an asymptotic explanation. Sections~10--11 supply generating functions,
sequence data, and the computational audit. The appendices collect the
polynomial formulas, the deliberate redundancies retained from the two
source drafts, a minimal symbolic check, and source-status details.
Throughout, watch the distinct roles of the row index $n$, the column
index $k$, the ground-set size $N=n+(r-1)k$, and the abbreviation
$t=n+(r-1)k-1$.
\end{tcolorbox}
\newpage

\section{The problem and the results}\label{sec:results}

Deb and Sokal introduce higher-order Stirling triangles and conjecture
row log-concavity through order five in
\cite[Conjectures~1.3(c) and 1.4(c)]{DebSokal}.
The cycle-number cases $r=3,4,5$ also appear as Problem~30.9 in the
BCC30 problem collection \cite{Cameron}. We select the subset-number
assertion because its recurrence has a particularly useful linear
coefficient. The second-order instance is the Ward triangle, recorded
as OEIS~A134991 after removing the zero column \cite{OEIS}.

Write $\assocS{N}{k}{r}$, also written $B_r(N,k)$ below, for the number of
set partitions of $[N]$ into $k$ blocks of size at least $r$, and
$\assocC{N}{k}{r}$ for the number of permutations of $[N]$ with $k$ cycles
of length at least $r$. Here $[N]=\{1,\ldots,N\}$ and the empty structure
has size zero. The arrays of interest are
\begin{equation}\label{eq:definitions}
 \Snum r n k=\assocS{n+(r-1)k}{k}{r},\qquad
 \Cnum r n k=\assocC{n+(r-1)k}{k}{r}.
\end{equation}
Their initial row is $(1,0,0,\ldots)$, and entries outside $0\le k\le n$
are zero. In every nonzero row, the positive entries are precisely those
with $1\le k\le n$.

\begin{definition}
A nonnegative sequence $(a_k)$ is \emph{log-concave} if
$a_k^2\ge a_{k-1}a_{k+1}$ at every relevant index. For a finite sequence
we extend it by zeros. We say it is \emph{strictly log-concave on the
interior of its positive support} when the inequality is strict whenever
all three entries are positive.
\end{definition}

\begin{theorem}[Sharp subset-number classification]\label{thm:main}
For a positive integer $r$, the following are equivalent:
\begin{enumerate}[label=(\roman*),nosep]
\item Every row of $S^{(r)}$ is log-concave.
\item $1\le r\le5$.
\end{enumerate}
For these five orders, and for every $n\ge2$ and $1\le k\le n-1$,
\begin{equation}\label{eq:main-strict}
 (\Snum r n k)^2>\Snum r n {k-1}\,\Snum r n {k+1}.
\end{equation}
For $r=6$ the first failing row is $n=4$: its relevant entries are
$(\Snum6 4 1,\Snum6 4 2,\Snum6 4 3)=(1,4719,23279256)$ and its
log-concavity defect is $-1010295$. For every $r\ge7$ the first failing
row is $n=3$, with the failure at $k=2$.
\end{theorem}

\begin{theorem}[Cycle-number consequence]\label{thm:cycles}
Every row of $C^{(r)}$ is log-concave for $1\le r\le4$, with strict
inequalities on the interior of the positive support. In particular,
the log-concavity cases $r=3,4$ of Deb--Sokal's Conjecture~1.3(c) hold,
as do the cases $r=3,4$ of BCC30 Problem~30.9.
\end{theorem}

These theorems concern the log-concavity clauses only. The same numbered
conjectures in \cite{DebSokal} contain other assertions; no claim about
all their clauses follows from this manuscript. In particular, we do not
settle the order-five cycle case, the general total positivity of the
triangles, or their Hankel-total-positivity questions. Nor do the local
inequalities proved here give total positivity of a row Toeplitz matrix:
\cref{ex:nonreal} exhibits a strictly log-concave row whose polynomial
has nonreal zeros, which already rules out the real-rootedness reading.

A crucial distinction is that $n$ in \eqref{eq:definitions} is \emph{not}
the cardinality of the underlying set. Moving right in a fixed row changes
that cardinality by $r-1$. Thus a theorem about a fixed-$N$ row of ordinary
associated Stirling numbers would not, by itself, solve this problem.
Likewise, these numbers must not be confused with the differently defined
$r$-Stirling numbers, which require certain distinguished elements to lie
in distinct blocks; compare the terminology in \cite[Section~1]{Shankar}.

\subsection{What makes the proof work}

The proof has three ingredients. First, a geometric rescaling removes the
factorial denominator from the subset recurrence. Second, three elementary
consequences of log-concavity reduce the next-row inequality to a quadratic
form in two adjacent entries of the previous row. Finally, the triangular
index restriction bounds its mixed coefficient. Orders three and four then
follow from nonnegative coefficients. At order five the mixed coefficient
can be negative --- at $n=3$, $k=2$ it equals $-4944$ --- so it must be
controlled rather than discarded. The control is exact: a single identity,
\eqref{eq:SOS}, writes twenty-five times the residual quadratic as a square
plus two terms whose nonnegativity is exactly the triangular domain
condition.

This is closely related to Sagan's recurrence criterion
\cite[Theorem~1]{Sagan}, whose nonnegative mixed-term hypothesis is sufficient
for some of the lower-order cases; the same coefficient-sign condition is
recorded as \cite[Theorem~1.2]{Shankar}. \dup{D8} Our local calculation
retains the positive square terms and allows a negative mixed term. We do
not claim that the general quadratic-form observation is historically new;
the specific application and factorization are what establish the selected
assertion here. Shankar's principal power-of-affine-coefficients family is
different from the falling-factorial weight used below.

\subsection{Provenance of this document}\label{subsec:provenance}

This manuscript merges two independently produced research packages that
proved the same theorem by the same four moves under different letters.
The first, in the repository directory
\texttt{docs/reports/log-concavity-and-unimodality/higher-order-stirling-rows},
was prepared on 19 September 2026; it proved the shared theorem by bounding
the mixed coefficient $H_d(t,k)$ below by its value at the non-integral
argument $k=t/(d+1)$, using strict convexity of the falling factorial
(\cref{lem:falling}), and then completing the square against the factored
discriminant $4A_4-L_4^2$. It additionally proved the cycle cases
$1\le r\le4$, the interpolating-family theorem, and the asymptotic
explanation of the cutoff. The second, in
\texttt{docs/reports/log-concavity-and-unimodality/stirling-subset-cutoff},
was prepared independently on 20 September 2026 and is recorded there as
having been drafted with ChatGPT; it proved the shared theorem by the exact
identity $H_4-L_4=\frac{96}{5}\beta(t)(t-5k)$, which replaces the convexity
argument by the immediate inequality $t-5k=n-k-1\ge0$, and folded the whole
interior step into the single nonnegative-summand identity~\eqref{eq:SOS}.

The present document keeps the first package as its spine and its notation,
takes the second package's identity as the primary fifth-order certificate,
and retains both routes wherever they constitute independent checks. The
twelve deliberate redundancies are tagged in the margin of the text as
\dup{D1}--\dup{D12} and are listed with their purpose in
\cref{sec:redundancy}.

\section{Recurrences, normalization, and boundary values}\label{sec:recurrence}

Set $d=r-1$, and define the falling-factorial polynomial
\begin{equation}\label{eq:P}
 P_d(t)=\prod_{j=0}^{d-1}(t-j),\qquad P_0(t)=1.
\end{equation}
The two recurrences in \cite[Lemma~1.1]{DebSokal} can be derived directly
from the distinguished largest label, as follows.

\begin{proposition}\label{prop:recurrences}
For $n\ge1$, $1\le k\le n$, and $t=n+dk-1$,
\begin{align}
 \Snum r n k
   &=\binom{t}{d}\Snum r {n-1}{k-1}+k\Snum r {n-1}k,
       \label{eq:Srec}\\
 \Cnum r n k
   &=P_d(t)\Cnum r {n-1}{k-1}+t\Cnum r {n-1}k.
       \label{eq:Crec}
\end{align}
\end{proposition}
\begin{proof}
Let $N=n+dk$. In a partition counted by $\assocS Nk r$, the block
containing $N$ either has exactly $r$ elements or has more than $r$.
In the first case choose its $r-1$ other elements and partition the remaining
$N-r$ elements into $k-1$ admissible blocks. In the second case remove $N$;
there is an admissible partition of $N-1$ elements and a choice of one of its
$k$ blocks in which to reinsert $N$. Hence
\[
 \assocS Nk r=\binom{N-1}{r-1}\assocS{N-r}{k-1}r
                 +k\assocS{N-1}kr.
\]
For permutations, a cycle of size exactly $r$ has $(r-1)!$ cyclic
orderings once its elements have been chosen. Otherwise, removing $N$
leaves a permutation into whose cycles $N$ can be reinserted after any
of the $N-1$ remaining labels. The corresponding weights are therefore
$(r-1)!\binom{N-1}{r-1}$ and $N-1$.
Finally,
$N-r=(n-1)+d(k-1)$ and $N-1=(n-1)+dk$, giving the stated recurrences.
\end{proof}

The argument includes $r=1$, where the first case means that $N$ is a
singleton. Since $B_r(N,k)$ vanishes for $0<N<rk$, the substitution
$N=n+(r-1)k$ makes the support exactly $0\le k\le n$: for $1\le k\le n$
one may take $k-1$ blocks of size $r$ and one block of size $r+n-k$, so
\begin{equation}\label{eq:support}
 \Snum r n0=0,\qquad \Snum r n1=1,\qquad \Snum r nk>0\quad(1\le k\le n)
\end{equation}
for $n\ge1$.

For the subset array define
\begin{equation}\label{eq:normalization}
 T^{(r)}_{n,k}=(d!)^k\Snum r n k.
\end{equation}
Then \eqref{eq:Srec} becomes
\begin{equation}\label{eq:Trec}
 T^{(r)}_{n,k}=P_d(n+dk-1)T^{(r)}_{n-1,k-1}
                      +kT^{(r)}_{n-1,k},
\end{equation}
an exact equivalence because $P_d(t)=d!\binom td$ and the factor $(d!)^k$
supplies one extra $d!$ in the first summand.
This normalization leaves log-concavity unchanged, since
\begin{equation}\label{eq:scale-difference}
 (T^{(r)}_{n,k})^2-T^{(r)}_{n,k-1}T^{(r)}_{n,k+1}
  =(d!)^{2k}\bigl[(\Snum r n k)^2-\Snum r n{k-1}\Snum r n{k+1}\bigr].
\end{equation}
The cycle array already has the same first weight $P_d(t)$; only its
second weight differs.

For later checks, direct counting gives, for $n\ge1$,
\begin{align}
 \Snum r n1&=1,&
 \Snum r nn&=\frac{(rn)!}{n!(r!)^n},\label{eq:Sboundary}\\
 \Cnum r n1&=(n+r-2)!,&
 \Cnum r nn&=\frac{(rn)!}{n!r^n}.\label{eq:Cboundary}
\end{align}
In particular $C^{(r)}_{1,1}=(r-1)!$; for $r>2$ this is not one.
These values also make the support assertions immediate from the recurrences.
The rows $n=0,1$ carry no internal inequality, and the only internal index
of row $n=2$ is $k=1$, handled in \cref{sec:subset-proof}.

\subsection{An independent counting formula}\label{subsec:profile}

A second, definition-level formula will be used for an independent check of
the small entries. A partition of $[N]$ with exactly $m_j$ blocks of size
$j$ is counted by
\begin{equation}\label{eq:profile}
 \frac{N!}{\prod_{j\ge r}(j!)^{m_j}\,m_j!},\qquad
 \sum_{j\ge r}m_j=k,\qquad\sum_{j\ge r}jm_j=N:
\end{equation}
order all $N$ elements, cut them into blocks of the prescribed sizes, and
divide out the order inside each block and the permutations of equally
sized blocks. Summing \eqref{eq:profile} over admissible block-size
profiles gives $B_r(N,k)$. This route uses neither the recurrence nor the
normalization, and it is different from the generating-function route of
\cref{sec:gf}; both are retained. \dup{D4}

\section{A sharp local quadratic-form criterion}\label{sec:criterion}

Let $(x_j)$ be a nonnegative log-concave sequence with no internal zeros.
At a fixed index $k$ suppose $x_{k-1}>0$ and $x_k>0$, and put
\[
 w=x_{k-2},\qquad u=x_{k-1},\qquad v=x_k,\qquad z=x_{k+1}.
\]

\begin{lemma}[Three slacks]\label{lem:cross}
With the preceding notation,
\begin{equation}\label{eq:three-slacks}
 wv\le u^2,\qquad uz\le v^2,\qquad wz\le uv.
\end{equation}
\end{lemma}
\begin{proof}
The first two are log-concavity at $k-1$ and $k$. For the third, if
$wz=0$ the claim is immediate; otherwise the interval-support condition
makes all four entries positive, and multiplying the first two
inequalities and dividing by $uv>0$ gives $wz\le uv$.
\end{proof}

Consider nonnegative coefficients and the local transform
\begin{equation}\label{eq:local-transform}
 y_j=a_jx_{j-1}+b_jx_j\qquad (j=k-1,k,k+1).
\end{equation}
Use the abbreviations $a=a_k$, $a_-=a_{k-1}$, $a_+=a_{k+1}$ and similarly
for $b$. Define
\begin{equation}\label{eq:ABH}
 A=a^2-a_-a_+,\qquad B=b^2-b_-b_+,\qquad
 H=2ab-a_-b_+-b_-a_+.
\end{equation}

\begin{lemma}[Exact slack decomposition]\label{lem:slacks}
With the preceding notation,
\begin{align}
 y_k^2-y_{k-1}y_{k+1}
  ={}& Au^2+Huv+Bv^2 \notag\\
    &+a_-a_+(u^2-wv)+b_-b_+(v^2-uz)
       +a_-b_+(uv-wz).\label{eq:slack-identity}
\end{align}
Consequently,
\begin{equation}\label{eq:quadratic-bound}
 y_k^2-y_{k-1}y_{k+1}\ge Au^2+Huv+Bv^2.
\end{equation}
\end{lemma}
\begin{proof}
Expand $(au+bv)^2-(a_-w+b_-u)(a_+v+b_+z)$ and collect terms.
This gives \eqref{eq:slack-identity} exactly. Each of the three remaining
brackets is nonnegative by \cref{lem:cross}, and each multiplier
is nonnegative.
\end{proof}

\begin{proposition}[Local preservation criterion]\label{prop:copositive}
The quadratic form $Au^2+Huv+Bv^2$ is nonnegative for every $u,v>0$
if and only if
\begin{equation}\label{eq:copositive}
 A\ge0,\qquad B\ge0,\qquad H\ge-2\sqrt{AB}.
\end{equation}
Thus these conditions imply the next-row log-concavity inequality.
They are also necessary if the transform is required to satisfy that
inequality for every positive log-concave four-term input.
\end{proposition}
\begin{proof}
If \eqref{eq:copositive} holds, write
\[
 Au^2+Huv+Bv^2=(\sqrt A\,u-\sqrt B\,v)^2
                         +(H+2\sqrt{AB})uv\ge0.
\]
Conversely, sending $v/u$ to zero or infinity forces $A\ge0$ and $B\ge0$.
For $A,B>0$, choosing $v/u=\sqrt{A/B}$ forces the third condition.
If $A=0$ or $B=0$, the same condition, now $H\ge0$, follows by taking
an appropriately small or large positive ratio.

For necessity as a preservation criterion, take the four input values
\[
 (w,u,v,z)=(u^2/v,\ u,\ v,\ v^2/u).
\]
They form a positive geometric progression, so all three slacks in
\eqref{eq:slack-identity} vanish. The next-row difference is exactly the
quadratic form. Extending this four-term sequence by zeros produces
no difficulty at the support boundaries.
\end{proof}

Only nonnegativity on the positive quadrant is needed, not positive
semidefiniteness on all of $\R^2$. In matrix terminology this is
copositivity of a $2\times2$ symmetric matrix. For exact computation the
square root in \eqref{eq:copositive} can be avoided entirely: require
$A,B\ge0$, and then require either $H\ge0$ or $H^2\le4AB$. For the hard
subset case we will use the stronger, more easily certified condition
$4AB-H^2>0$ after replacing $H$ by a suitable lower bound --- and in fact
even that check is replaced, in \cref{lem:certificate}, by an identity
with visibly nonnegative terms.

\begin{remark}[Why a negative mixed term is not fatal]\label{rem:sagan}
Sagan's sufficient condition \cite[Theorem~1]{Sagan}, and equally the
coefficient-sign criterion \cite[Theorem~1.2]{Shankar}, correspond to
$A,B,H\ge0$. If $H<0$, rejecting the recurrence at once loses the positive
contributions $Au^2$ and $Bv^2$. The exact square-completion threshold in
\eqref{eq:copositive} is what allows us to handle the fifth-order subset
array and the fourth-order cycle array. \dup{D8}
\end{remark}

\section{Proof of subset log-concavity through order five}\label{sec:subset-proof}

Assume row $n-1$ of $T^{(r)}$ is log-concave. The index $k=1$ is immediate
from \eqref{eq:support} and \eqref{eq:scale-difference}: since
$\Snum r n0=0$ and $\Snum r n1=1$, the normalized defect at $k=1$ equals
$(d!)^2>0$. For the remaining internal indices fix
\begin{equation}\label{eq:domain}
 n\ge3,\qquad 2\le k\le n-1,\qquad d=r-1,\qquad t=n+dk-1.
\end{equation}
For \eqref{eq:Trec}, the local weights are
\[
 a=P_d(t),\quad a_-=P_d(t-d),\quad a_+=P_d(t+d),\qquad
 b=k,\quad b_-=k-1,\quad b_+=k+1.
\]
Notice that the step in $t$ between \emph{adjacent columns of one row} is
$d$, not $d+1$. The latter number enters a different way: since $n\ge k+1$,
\begin{equation}\label{eq:triangular-bound}
 t=(d+1)k+(n-k-1)\ge(d+1)k,\qquad k\le\frac{t}{d+1},\qquad t\ge2(d+1).
\end{equation}
Treating $t$ and $k$ as unconstrained independent parameters would discard
exactly the information the fifth-order argument needs.

Let
\begin{align}
 A_d(t)&=P_d(t)^2-P_d(t-d)P_d(t+d),\label{eq:Ad}\\
 H_d(t,k)&=2kP_d(t)-(k+1)P_d(t-d)-(k-1)P_d(t+d).
                                                       \label{eq:Hd}
\end{align}
The remaining quadratic coefficient is simply $B=1$.

\subsection{Two elementary properties of the falling factorial}

\begin{lemma}\label{lem:falling}
For $d\ge1$ and $t>2d-1$, $A_d(t)>0$. For $d\ge2$ in the same domain,
\[
 \Delta_d(t):=P_d(t-d)+P_d(t+d)-2P_d(t)>0.
\]
For $d=0,1$, the centered second difference is zero.
\end{lemma}
\begin{proof}
For $d\ge1$,
\[
 P_d(t-d)P_d(t+d)=\prod_{j=0}^{d-1}\bigl((t-j)^2-d^2\bigr).
\]
Every factor is positive and strictly smaller than $(t-j)^2$, because
$t>2d-1$. Comparing products proves $A_d(t)>0$.

For $d\ge2$, $P_d''(s)>0$ whenever $s>d-1$: differentiate the product
twice, obtaining a sum of positive products of the remaining $d-2$
factors. Thus $P_d$ is strictly convex on the interval $[t-d,t+d]$,
and its centered second difference is positive. The polynomials
$P_0=1$ and $P_1=t$ have zero centered second difference.
\end{proof}

This lemma is used three times: for $A_d>0$ in every order treated here,
for the interpolating family of \cref{thm:interpolation}, and for the
asymptotics of \cref{sec:asymptotics}. Rewriting the mixed coefficient as
\begin{equation}\label{eq:Hdecreasing}
 H_d(t,k)=P_d(t+d)-P_d(t-d)-k\Delta_d(t)
\end{equation}
shows that for fixed $t$ it is affine and nonincreasing in $k$. Hence
\eqref{eq:triangular-bound} yields the lower bound
\begin{equation}\label{eq:Lbound}
 H_d(t,k)\ge L_d(t):=H_d\!\left(t,\frac{t}{d+1}\right).
\end{equation}
The argument $t/(d+1)$ need not be an integer; $H_d(t,k)$ is affine in $k$,
so this substitution is legitimate.

\begin{remark}[Two routes to the mixed-coefficient bound]\label{rem:tworoutes}
At $d=4$ the bound \eqref{eq:Lbound} is not needed in this form: the exact
identity \eqref{eq:MminusL} below gives $H_4-L_4=\frac{96}{5}\beta(t)(t-5k)$,
so $H_4\ge L_4$ follows from $t-5k=n-k-1\ge0$ and $\beta>0$ with no appeal
to convexity at all. That identity is the route used in the proof. The
convexity route is retained because it is the only one of the two that
survives to general $d$: \cref{thm:interpolation} and
\cref{sec:asymptotics} both depend on $\Delta_d(t)>0$. \dup{D2}
\end{remark}

\subsection{The orders $r=1,2,3,4$}

Direct expansion gives
\begin{center}
\renewcommand{\arraystretch}{1.45}
\begin{tabular}{ccl}
\toprule
$r$ & $d$ & $L_d(t)$\\
\midrule
1 & 0 & $0$\\
2 & 1 & $2$\\
3 & 2 & $\dfrac43(4t-3)$\\
4 & 3 & $\dfrac32(3t^2-15t+44)$\\
\bottomrule
\end{tabular}
\end{center}
For $d=0$, $A_0=0$, $H_0=0$, and the quadratic lower bound is $v^2>0$.
For $d=1$, $A_1=1$ and $H_1=2$, so it is $(u+v)^2>0$.
For $d=2$, $H_2=4(2t-2k-1)$ and $k\le t/3$ give $H_2\ge\frac43(4t-3)>0$
because $t\ge6$.
For $d=3$, positivity follows even without the domain restriction from
\[
 3t^2-15t+44=3(t-5/2)^2+101/4>0,
\]
equivalently from the negative discriminant $225-528<0$.
In the latter two cases $A_d>0$ by \cref{lem:falling}, so the quadratic
form is strictly positive for $u,v>0$.

\subsection{The order $r=5$}

Here $d=4$, $t=n+4k-1\ge5k\ge10$, and
\begin{align}
 P_4(t)&=t(t-1)(t-2)(t-3),\label{eq:P4}\\
 A_4(t)&=32\bigl(2t^6-18t^5+17t^4+168t^3
                      -55t^2-834t-630\bigr),\label{eq:A4}\\
 H_4(t,k)&=16\bigl(-12kt^2+36kt-54k+2t^3-9t^2+43t-51\bigr),
                                                       \label{eq:H4}\\
 L_4(t)&=-\frac{16}{5}\bigl(2t^3+9t^2-161t+255\bigr).
                                                       \label{eq:L4}
\end{align}
Unlike the lower orders, $H_4$ is not always nonnegative: at $n=3$, $k=2$,
so $t=10$, it equals $-4944$. The coefficient-sign criterion of
\cref{rem:sagan} therefore cannot prove this case, and the quadratic form
must be used in full.

Introduce the three small polynomials
\begin{equation}\label{eq:abc}
 \alpha(t)=2t^3+9t^2-161t+255,\quad
 \beta(t)=2t^2-6t+9,\quad
 \gamma(t)=7t^3-31t^2-126t+1080,
\end{equation}
so that $L_4(t)=-\frac{16}{5}\alpha(t)$. (These are the polynomials written
$A,B,C$ in the second source draft; they are renamed here to avoid a
collision with the coefficients $A_d$ and $B$ of \eqref{eq:ABH}.)

\begin{lemma}[Integer square certificate]\label{lem:certificate}
For every real $t,k,u,v$,
\begin{equation}\label{eq:SOS}
\boxed{\begin{aligned}
 25\bigl(A_4(t)u^2+H_4(t,k)uv+v^2\bigr)
  ={}&\bigl(5v-8\alpha(t)u\bigr)^2\\
    &+96(t-5)\beta(t)\gamma(t)u^2\\
    &+480\,\beta(t)\,(t-5k)\,uv.
\end{aligned}}
\end{equation}
If $t\ge10$, $t\ge5k$, and $u,v\ge0$, the right-hand side is nonnegative,
and it is strictly positive whenever $u>0$.
\end{lemma}
\begin{proof}
The identity is verified by multiplying out both sides. A less bulky
verification uses the two exact identities
\begin{align}
 H_4(t,k)-L_4(t)&=\frac{96}{5}\beta(t)\,(t-5k),\label{eq:MminusL}\\
 4A_4(t)-L_4(t)^2&=\frac{384}{25}(t-5)\beta(t)\gamma(t).
                                             \label{eq:main-factor}
\end{align}
Indeed, \eqref{eq:Hdecreasing} shows that $H_4$ is affine in $k$ with slope
$-\Delta_4(t)=-96\beta(t)$ and that $L_4=H_4(t,t/5)$, which is exactly
\eqref{eq:MminusL}. Completing the square in $v$ gives
\[
 A_4u^2+L_4uv+v^2=\left(v+\frac{L_4}2u\right)^2
   +\frac{4A_4-L_4^2}{4}u^2 ,
\]
and multiplying by $25$, using $\frac52L_4=-8\alpha$ and
\eqref{eq:main-factor}, turns this into the first two summands of
\eqref{eq:SOS}. Adding $25(H_4-L_4)uv$ and applying \eqref{eq:MminusL}
supplies the third. Equivalently, \eqref{eq:SOS} is the conjunction of the
three coefficientwise identities
\[
 25A_4=64\alpha^2+96(t-5)\beta\gamma,\qquad
 25H_4=-80\alpha+480\beta\,(t-5k),\qquad 25=25 ,
\]
each of which is a polynomial identity in $t$ alone once the affine
$k$-dependence is separated. All of them are checked by exact coefficient
comparison in the accompanying \texttt{certificates.py}, which uses no
computer algebra system, and again in \texttt{code/verify\_certificates.py},
which uses SymPy. \dup{D5}

For the signs, first
\begin{equation}\label{eq:Bpositive}
 \beta(t)=2\left(t-\tfrac32\right)^2+\tfrac92>0
\end{equation}
for all real $t$, and, writing $t=10+s$ with $s\ge0$,
\begin{equation}\label{eq:positive-cubic}
 \gamma(10+s)=7s^3+179s^2+1354s+3720>0 .
\end{equation}
Also $t-5>0$ and $t-5k\ge0$ on the domain. Hence all three summands of
\eqref{eq:SOS} are nonnegative for $u,v\ge0$, and the middle one is
strictly positive when $u>0$.
\end{proof}

Two further forms of the same certificate are kept, because each is
checkable in a way the others are not. \dup{D1} The two-step form is the
factored discriminant \eqref{eq:main-factor} together with the completed
square
\begin{align}
 A_4(t)u^2+H_4(t,k)uv+v^2
 &\ge A_4(t)u^2+L_4(t)uv+v^2\notag\\
 &=\left(v+\frac{L_4(t)}2u\right)^2
       +\frac{4A_4(t)-L_4(t)^2}{4}\,u^2
 >0;\label{eq:hard-square}
\end{align}
this is the form that generalizes, unchanged in shape, to the cycle
coefficients of \cref{sec:cycle-proof} and to arbitrary $d$ in
\cref{sec:asymptotics}. The fully expanded coefficient form of
\eqref{eq:main-factor} is recorded in \cref{sec:polynomial-table}; it is
the only version that a reader can check entirely by hand.

\subsection{The interior induction step in one display}

Collecting \cref{lem:slacks} and \cref{lem:certificate} gives a single
equation that is the whole interior step of the induction. \dup{D1}

\begin{proposition}[Six-term defect certificate]\label{prop:full}
Let $r=5$, let $n\ge3$ and $2\le k\le n-1$, and put
\[
 a_j=T^{(5)}_{n-1,j},\quad b_j=T^{(5)}_{n,j},\quad
 t=n+4k-1,\quad u=a_{k-1},\quad v=a_k .
\]
Then, with $P_4,\alpha,\beta,\gamma$ as above,
\begin{align}
 25\bigl(b_k^2-b_{k-1}b_{k+1}\bigr)
 ={}&25\,P_4(t-4)P_4(t+4)\bigl(a_{k-1}^2-a_{k-2}a_k\bigr)\notag\\
 &+25(k^2-1)\bigl(a_k^2-a_{k-1}a_{k+1}\bigr)\notag\\
 &+25\,P_4(t-4)(k+1)\bigl(a_{k-1}a_k-a_{k-2}a_{k+1}\bigr)\notag\\
 &+\bigl(5v-8\alpha(t)u\bigr)^2\notag\\
 &+96(t-5)\beta(t)\gamma(t)\,u^2\notag\\
 &+480\,\beta(t)\,(n-k-1)\,uv.\label{eq:full-certificate}
\end{align}
\end{proposition}
\begin{proof}
Multiply \eqref{eq:slack-identity} by $25$, with the weights
$a_-a_+=P_4(t-4)P_4(t+4)$, $b_-b_+=(k-1)(k+1)$ and
$a_-b_+=P_4(t-4)(k+1)$, and replace the resulting
$25(A_4u^2+H_4uv+v^2)$ by the right-hand side of \eqref{eq:SOS}, using
$t-5k=n-k-1$.
\end{proof}

The signs of the six summands require, in order: log-concavity of row
$n-1$ at $k-1$; log-concavity at $k$; \cref{lem:cross}; that a square is
nonnegative; $t\ge10$ together with \eqref{eq:Bpositive} and
\eqref{eq:positive-cubic}; and $n-k-1\ge0$. The fifth summand is strictly
positive because $u>0$.

\subsection{Completing the induction, including the edges}

The rows $n=0,1$ have no internal inequality, and row $n=2$ has only the
index $k=1$ treated above. Assume row $n-1$ is log-concave.
For $2\le k\le n-1$, both $u=T^{(r)}_{n-1,k-1}$ and
$v=T^{(r)}_{n-1,k}$ are positive. At $k=2$, $w=0$; at $k=n-1$, $z=0$.
Neither zero causes a problem in \cref{lem:slacks}. The preceding
calculations show that its quadratic lower bound is strictly positive for
every order $r\le5$; for $r=5$ this is \cref{lem:certificate}, or in one
line \eqref{eq:full-certificate}. Hence all the positive-support interior
inequalities in row $n$ are strict. At either endpoint of the positive
support one neighboring entry is zero, and outside the support the
inequalities are trivial. This completes the induction for $T^{(r)}$.
Equation \eqref{eq:scale-difference} transfers the result to $S^{(r)}$.

Three features of this induction deserve emphasis. Nothing in the argument
requires strict log-concavity of the input row: ordinary log-concavity, the
interval support condition, and $u>0$ already suffice, so the induction
\emph{creates} strictness rather than silently assuming it. The last
internal index $k=n-1$ is genuinely included, since there $a_{k+1}=0$ and
$t-5k=0$, both of which the certificate allows. And the smallest genuine
case $n=3$, $k=2$, $t=10$ is covered without an asymptotic exception.

We have proved the sufficiency direction of \cref{thm:main}, including the
strict inequalities, without using real-rootedness or a finite
verification bound.

\section{An explicit fifth-order \Turan{} inequality}\label{sec:quantitative}

The certificate gives more than the sign of a difference. Let
\begin{equation}\label{eq:Theta}
 \Theta(t)=(t-5)\beta(t)\gamma(t)
   =(t-5)(2t^2-6t+9)(7t^3-31t^2-126t+1080).
\end{equation}

\begin{corollary}\label{cor:quantitative}
For $n\ge3$, $2\le k\le n-1$, and $t=n+4k-1$,
\begin{equation}\label{eq:quantitative}
 (\Snum5nk)^2-\Snum5n{k-1}\Snum5n{k+1}
 \ge\frac{\Theta(t)}{150}\,(\Snum5{n-1}{k-1})^2>0.
\end{equation}
\end{corollary}
\begin{proof}
Drop every summand of \eqref{eq:full-certificate} except the fifth.
The difference for the normalized row is then at least
$(96/25)\Theta(t)(T^{(5)}_{n-1,k-1})^2$.
Now divide by $24^{2k}$ and use
$T^{(5)}_{n-1,k-1}=24^{k-1}S^{(5)}_{n-1,k-1}$.
The remaining constant is $96/(25\cdot24^2)=1/150$.
\end{proof}

For instance, the nonzero portion of row three is
\[
 (S^{(5)}_{3,1},S^{(5)}_{3,2},S^{(5)}_{3,3})=(1,462,126126).
\]
Its middle difference is $462^2-126126=87318$.
The lower bound in \eqref{eq:quantitative}, with $t=10$, is $18476$.
It is not asserted to be optimal; its point is to make strict positivity
uniform and explicit. At that same smallest index the scalar coefficients
are
\[
 A_4(10)=16752960,\quad H_4(10,2)=-4944,\quad
 \alpha(10)=1545,\quad\beta(10)=149,\quad\gamma(10)=3720,
\]
so that $4A_4-H_4^2=42568704>0$ although $H_4<0$: this audits the smallest
nontrivial fifth-order case and shows concretely how much the positive
square terms retain.

\subsection{Shape consequences}

In a positive row, \eqref{eq:main-strict} is equivalent to strict decrease
of successive ratios:
\begin{equation}\label{eq:ratio-monotonicity}
 \frac{S^{(r)}_{n,2}}{S^{(r)}_{n,1}}>
 \frac{S^{(r)}_{n,3}}{S^{(r)}_{n,2}}>\cdots>
 \frac{S^{(r)}_{n,n}}{S^{(r)}_{n,n-1}} .
\end{equation}
Therefore every such row increases until its successive ratios cross $1$
and then decreases; it is unimodal, with either a unique maximum or
two adjacent equal maxima, since strict log-concavity does not exclude the
latter. Multiplying entries by $q^k$, for any $q>0$, preserves the theorem.
Normalizing these weighted entries by their sum gives a log-concave
probability mass function on the column index. This is a distribution over
a row of the shifted array, not over partitions of a single fixed set size.

The same ratio comparison yields a formally stronger statement, which does
not follow from the unimodality remark and does not imply it: \dup{D11}
\begin{equation}\label{eq:balanced}
 S^{(r)}_{n,i}S^{(r)}_{n,j}>S^{(r)}_{n,i-1}S^{(r)}_{n,j+1}
 \qquad(2\le i\le j\le n-1).
\end{equation}
Indeed, compare the ratio at $i$ with the ratio at $j+1$ in
\eqref{eq:ratio-monotonicity} and multiply by the positive denominators.
Cases involving a zero outside the support satisfy the corresponding weak
inequality. No separate combinatorial injection is needed.

\subsection{Log-concavity without real-rootedness}

\begin{example}\label{ex:nonreal}
For $r=5$ and $n=3$ the row polynomial is
\[
 x+462x^2+126126x^3=x\bigl(1+462x+126126x^2\bigr).
\]
The row is strictly log-concave, since $462^2-126126=87318>0$, yet the
quadratic factor has discriminant
\[
 462^2-4\cdot126126=-291060<0,
\]
so its roots are not real.
\end{example}

This example is a scope fence. \dup{D12} Any argument that tried to prove
real-rootedness first would be attempting a statement that is false for
this family. It also shows that the local inequalities established here
cannot yield total positivity of the triangle or nonnegativity of all
minors of a row Toeplitz matrix. The complementary fence is the
result-to-source map in \cref{sec:status}, which records exactly which
clause of which conjecture is and is not settled.

\section{Why order five is the exact cutoff}\label{sec:sharpness}

\subsection{The sixth-order counterexample}

For $r=6$ the first relevant rows are
\begin{align*}
 (S^{(6)}_{2,k})_{k=0}^{2}&=(0,1,462),\\
 (S^{(6)}_{3,k})_{k=0}^{3}&=(0,1,1716,2858856),\\
 (S^{(6)}_{4,k})_{k=0}^{4}&=(0,1,4719,23279256,96197645544).
\end{align*}
Row three passes, because $1716^2-2858856=85800>0$. The two entries that
decide row four can be obtained in two independent ways, and both are
recorded here because this is the single load-bearing numerical fact of
the whole converse. \dup{D3} By block-size pattern,
\begin{align*}
 S^{(6)}_{4,2}&=\binom{14}{6}+\frac12\binom{14}{7}=4719,\\
 S^{(6)}_{4,3}&=\frac{19!}{2!\,(6!)^2\,7!}=23279256,
\end{align*}
the second entry counting the patterns $(6,8)$ and $(7,7)$ and the third
the pattern $(6,6,7)$. By the recurrence \eqref{eq:Srec},
\begin{align*}
 S^{(6)}_{4,2}&=\binom{13}{5}+2\cdot1716=4719,\\
 S^{(6)}_{4,3}&=\binom{18}{5}\cdot1716+3\cdot2858856=23279256 .
\end{align*}
Consequently,
\begin{equation}\label{eq:r6-counterexample}
 (S^{(6)}_{4,2})^2-S^{(6)}_{4,1}S^{(6)}_{4,3}
 =22268961-23279256=-1010295<0.
\end{equation}
Rows zero and one have no internal test, and the only test for row two has
a zero left neighbour. Together with the row-three calculation, this proves
that row four is the first failure for $r=6$.

\subsection{A row-three counterexample for every $r\ge7$}

The nonzero part of row three has a uniform closed form:
\begin{equation}\label{eq:row3}
 S^{(r)}_{3,1}=1,\qquad
 S^{(r)}_{3,2}=\binom{2r+1}{r},\qquad
 S^{(r)}_{3,3}=\frac{(3r)!}{6(r!)^3}.
\end{equation}
The two-block case has block sizes $r,r+1$, and the three-block case
has three blocks of size $r$, as in \eqref{eq:Sboundary}.
Define the ratio that decides log-concavity in this row:
\begin{equation}\label{eq:Fr}
 F(r)=\frac{\binom{2r+1}{r}^{\!2}}{(3r)!/[6(r!)^3]},\qquad
 R(r)=\frac1{F(r)}=\frac{(3r)!\,(r+1)!^2}{6\,r!\,(2r+1)!^2}.
\end{equation}
The row-three test fails exactly when $F(r)<1$, equivalently when $R(r)>1$.
Both presentations are kept, since each is a transcription check on the
other. \dup{D9} Cancellation of consecutive factorials gives
\begin{equation}\label{eq:F-ratio}
 \frac{F(r+1)}{F(r)}
  =\frac{4(r+1)^2(2r+3)^2}{3(r+2)^2(3r+2)(3r+1)},\qquad
 \frac{R(r+1)}{R(r)}
  =\frac{3(3r+1)(3r+2)(r+2)^2}{4(r+1)^2(2r+3)^2}.
\end{equation}
In either form, the denominator of $F(r+1)/F(r)$ minus its numerator is
\begin{equation}\label{eq:ratio-gap}
 11r^4+55r^3+74r^2+12r-12>0\qquad(r\ge1).
\end{equation}
For a coefficientwise positivity certificate, substitute $r=1+w$; the
polynomial becomes
\[
 11w^4+99w^3+305w^2+369w+140 .
\]
Thus $F(r)$ is strictly decreasing, and $R(r)$ strictly increasing, on the
positive integers. At $r=7$,
\[
 \binom{15}{7}^2=6435^2=41409225<66512160
             =\frac{21!}{6(7!)^3}.
\]
Therefore $F(r)<1$ for every $r\ge7$, proving a failure in row three.
No earlier row can fail because it has at most two positive entries.
Together with \eqref{eq:r6-counterexample}, this proves necessity and
completes \cref{thm:main}.

\subsection{The third row at the threshold}

\begin{table}[htbp]
\centering
\small
\renewcommand{\arraystretch}{1.2}
\begin{tabular}{rrrr}
\toprule
$r$ & $S^{(r)}_{3,2}$ & $S^{(r)}_{3,3}$ &
$(S^{(r)}_{3,2})^2-S^{(r)}_{3,3}$\\
\midrule
1 & 3 & 1 & 8\\
2 & 10 & 15 & 85\\
3 & 35 & 280 & 945\\
4 & 126 & 5775 & 10101\\
5 & 462 & 126126 & 87318\\
6 & 1716 & 2858856 & 85800\\
7 & 6435 & 66512160 & $-25102935$\\
8 & 24310 & 1577585295 & $-986609195$\\
\bottomrule
\end{tabular}
\caption{Exact third-row data; $S^{(r)}_{3,1}=1$ for every $r$.
The order-six line is the reason the converse needs row four.}
\label{tab:third}
\end{table}

\cref{tab:third} makes the structure of the converse visible. The
order-six row three still passes, and only its row four fails. The
higher-order obstruction is therefore not a uniform single-row phenomenon
beginning at order six: an argument that checked only row three would
locate the cutoff at the wrong place.

\section{Cycle triangles and an interpolating family}\label{sec:cycle-proof}

For the cycle recurrence, the first quadratic coefficient is the same
$A_d(t)$, but
\[
 b=t,\quad b_-=t-d,\quad b_+=t+d,\quad B=d^2.
\]
The mixed coefficient is
\begin{equation}\label{eq:Md}
 M_d(t)=2tP_d(t)-(t+d)P_d(t-d)-(t-d)P_d(t+d).
\end{equation}
For $d=1,2,3$ it is
\begin{align}
 M_1(t)&=2,\label{eq:M1}\\
 M_2(t)&=8(t-1),\label{eq:M2}\\
 M_3(t)&=-18(3t-11).\label{eq:M3}
\end{align}
The first two cases have a positive mixed coefficient. For $d=3$, use
\begin{equation}\label{eq:cycle4-discriminant}
 4\cdot9A_3(t)-M_3(t)^2
      =972(t-3)^2(t-1)(t+3)>0\qquad(t\ge8).
\end{equation}
Since $A_3(t)>0$ and $B=9$, the quadratic form is positive definite.
This proves strict log-concavity for orders $r=2,3,4$ by the same induction
and boundary argument as in \cref{sec:subset-proof}. Note that the proof
here is the completed-square form \eqref{eq:hard-square}, not the
single-identity form \eqref{eq:SOS}: the two-step route is what transfers
to a different second weight.

For completeness, ordinary cycle numbers ($r=1$, $d=0$) can also be
handled without a real-rootedness theorem. The recurrence has $a=1$ and
$b=n-1$, independent of $k$, so $A=B=M=0$. The first slack in
\eqref{eq:slack-identity} nevertheless supplies strictness. For $k=2$,
$w=0$ and $u>0$, hence $u^2-wv>0$. For $k>2$, strict log-concavity in
row $n-1$ at index $k-1$ gives the same inequality. The base row with two
positive entries is automatic. This completes \cref{thm:cycles}.

\subsection{A continuous extension}

The subset and cycle results of orders at most four fit into a single
family. The coefficients below need not be integers.

\begin{theorem}\label{thm:interpolation}
Let $1\le d\le3$, let $\rho,\sigma\ge0$ with $\rho+\sigma>0$, and define
an array $U$ with initial row $U_{0,k}=\delta_{0k}$, zero entries outside
$0\le k\le n$, and recurrence
\begin{equation}\label{eq:interpolating}
 U_{n,k}=P_d(n+dk-1)U_{n-1,k-1}
       +\bigl[\rho(n+dk-1)+\sigma k\bigr]U_{n-1,k}
\end{equation}
for $1\le k\le n$, with $U_{n,0}=0$ for $n\ge1$.
Every row is strictly log-concave on the interior of its positive support.
\end{theorem}
\begin{proof}
All entries on $1\le k\le n$ are positive. At a fixed interior index,
write $t=n+dk-1$. The first quadratic coefficient remains $A_d(t)$.
Since the second weight is affine in $k$ with slope $\rho d+\sigma$,
\[
 B=(\rho d+\sigma)^2,\qquad
 H=\rho M_d(t)+\sigma H_d(t,k).
\]
Our preceding calculations give
\[
 M_d(t)>-2d\sqrt{A_d(t)},\qquad
 H_d(t,k)>-2\sqrt{A_d(t)}
\]
for $d=1,2,3$ in the triangular interior; the second bound uses
\eqref{eq:Lbound}, hence \cref{lem:falling}. Multiplying by $\rho,\sigma$
and adding yields
\[
 H>-2(\rho d+\sigma)\sqrt{A_d(t)}=-2\sqrt{A_d(t)B}.
\]
The local criterion proves strict positivity of the quadratic form, and
induction completes the proof.
\end{proof}

The choice $(\rho,\sigma)=(0,1)$ is the normalized subset array, and
$(1,0)$ is the cycle array. This theorem is a genuine parameter extension,
not a convex-combination assertion about already computed rows: the
parameters are inserted in the recurrence at every step.

\section{A rigorous obstruction in the fifth-order cycle case}\label{sec:obstruction}

The preceding method does not automatically settle the remaining cycle
case. In fact, universal preservation of log-concavity is false for its
row operator, and we record two independent witnesses of this. \dup{D7}
At $d=4$,
\begin{equation}\label{eq:M4}
 M_4(t)=-32(t-2)(2t^2+4t-51),
\end{equation}
and the algebraic witness is the factorization
\begin{equation}\label{eq:cycle5-discriminant}
 \boxed{64A_4(t)-M_4(t)^2
       =-9216(t-4)(t+4)(2t-9)\beta(t).}
\end{equation}
For $t\ge10$ the right-hand side is negative and $M_4(t)<0$.
Thus the necessary condition of \cref{prop:copositive} fails.

The numerical witness is an exact input that also respects the support of a
preceding triangular row. Let $n=5$, $q=7350$, and take
\[
 (x_0,x_1,x_2,x_3,x_4,x_5)=(0,1,q,q^2,q^3,0).
\]
It is log-concave. Apply the fifth-order cycle row operator
\[
 y_j=P_4(4+4j)x_{j-1}+(4+4j)x_j\qquad(1\le j\le5).
\]
The three relevant output entries are
\begin{equation}\label{eq:witness-values}
 y_2=100080,\qquad y_3=1185408000,\qquad y_4=14223043800000,
\end{equation}
and their \Turan{} difference is
\begin{equation}\label{eq:witness-negative}
 y_3^2-y_2y_4=-18250097040000000<0.
\end{equation}
The choice of $q$ is not arbitrary. At this index $t=16$, the quadratic
form $A_4u^2+M_4uv+16v^2$ is minimized as a function of $v/u$ at
\[
 \frac vu=-\frac{M_4(16)}{32}
       =(16-2)\bigl(2\cdot16^2+4\cdot16-51\bigr)=7350.
\]
The geometric input makes every slack zero and realizes the negative
quadratic value exactly.

\begin{tcolorbox}
\textbf{What this does and does not show.}
The artificial input is not a row of $C^{(5)}$.
Equation~\eqref{eq:witness-negative}, like the factorization
\eqref{eq:cycle5-discriminant}, disproves only the claim that this
operator preserves \emph{all} log-concave inputs. Rows generated by one
particular recurrence can satisfy additional inequalities that arbitrary
log-concave inputs do not. Neither statement disproves log-concavity of the
actual fifth-order Stirling cycle rows, and no resolution of that case ---
Deb--Sokal Conjecture~1.3(c) for $r=5$, equivalently BCC30 Problem~30.9
for $r=5$ --- is claimed here. What \eqref{eq:cycle5-discriminant} locates
is the exact point at which the present universal-preservation argument
ceases to apply.
\end{tcolorbox}

There is a concrete next invariant to investigate. For the actual previous
row, retain the three nonnegative slacks in \eqref{eq:slack-identity}
and seek a lower bound on their weighted sum that dominates the possible
negative part of $A_4u^2+M_4uv+16v^2$. Alternatively, prove a restriction
on the actual adjacent ratio $v/u$ that excludes the dangerous interval
between the two positive roots of
\[
 16q^2+M_4(t)q+A_4(t)=0.
\]
Either approach would use information specific to the array, which
ordinary row log-concavity alone cannot supply. These are proposed
extensions, not additional proved statements.

\section{A leading-order explanation of the cutoff}\label{sec:asymptotics}

The exact proof is polynomial, but its structure has a useful asymptotic
interpretation. Fix $d\ge1$ and let $t\to\infty$. Since
$P_d(t)=t^d+O(t^{d-1})$, centered Taylor expansions give
\begin{align}
 P_d(t+d)-P_d(t-d)&=2d^2t^{d-1}+O(t^{d-2}),\label{eq:asym-diff}\\
 \Delta_d(t)&=d^3(d-1)t^{d-2}+O(t^{d-3}).\label{eq:asym-second}
\end{align}
For $d=1$ the centered second difference is exactly zero, as this formula
permits. Comparing the leading terms in $P_d^2-P_d(t-d)P_d(t+d)$ gives
\begin{equation}\label{eq:asym-A}
 A_d(t)=d^3t^{2d-2}+O(t^{2d-3}).
\end{equation}
For example, this also follows from
$\log P_d(t-d)+\log P_d(t+d)-2\log P_d(t)
=-d^3/t^2+O(t^{-3})$.
Combining \eqref{eq:asym-diff}--\eqref{eq:asym-second} with
$L_d=P_d(t+d)-P_d(t-d)-t\Delta_d/(d+1)$ yields
\begin{equation}\label{eq:asym-L}
 L_d(t)=\frac{d^2(-d^2+3d+2)}{d+1}\,t^{d-1}+O(t^{d-2}).
\end{equation}
Hence the normalized mixed coefficient has limit
\begin{equation}\label{eq:asym-subset-ratio}
 \frac{L_d(t)}{2\sqrt{A_d(t)}}
 \longrightarrow \frac{\sqrt d\,(-d^2+3d+2)}{2(d+1)}.
\end{equation}
At $d=4$ this is $-2/5$, safely above the local threshold $-1$.
At $d=5$ it is $-2\sqrt5/3<-1$. Thus the elementary preservation
mechanism naturally changes between orders five and six. The exact
counterexamples in \cref{sec:sharpness}, rather than this asymptotic
observation alone, prove the cutoff for the actual arrays.

For cycles, \eqref{eq:Md} can be written as
$M_d=d(P_d(t+d)-P_d(t-d))-t\Delta_d$. Therefore
\begin{equation}\label{eq:asym-M}
 M_d(t)=d^3(3-d)t^{d-1}+O(t^{d-2}),\qquad
 \frac{M_d(t)}{2d\sqrt{A_d(t)}}\longrightarrow
                         \frac{\sqrt d(3-d)}2.
\end{equation}
At $d=4$ the latter limit is exactly $-1$, on the boundary.
The leading terms of $64A_4$ and $M_4^2$ cancel; their difference begins
with $-36864t^5$. This explains why the fifth-order cycle operator is
more delicate than its subset counterpart, and why keeping lower-order
terms is essential.

\section{Generating functions and reproducible sequence data}\label{sec:gf}

\subsection{An independent coefficient formula}

Define formal power series
\begin{equation}\label{eq:fg}
 f_r(z)=\sum_{j\ge0}\frac{z^j}{(r+j)!},\qquad
 g_r(z)=\sum_{j\ge0}\frac{z^j}{r+j}.
\end{equation}
For each fixed $k$, counting ordered collections of $k$ blocks and then
dividing by $k!$ gives
\begin{align}
 \sum_{N\ge0}\assocS Nkr\frac{z^N}{N!}
  &=\frac1{k!}\left(\sum_{m\ge r}\frac{z^m}{m!}\right)^k,
                                                       \label{eq:associated-S-egf}\\
 \sum_{N\ge0}\assocC Nkr\frac{z^N}{N!}
  &=\frac1{k!}\left(\sum_{m\ge r}\frac{z^m}{m}\right)^k.
                                                       \label{eq:associated-C-egf}
\end{align}
For cycles, the coefficient $1/m$ is $(m-1)!/m!$.
Substituting $N=n+(r-1)k$ and factoring out $z^{rk}$ gives
\begin{align}
 \Snum rnk&=\frac{(n+(r-1)k)!}{k!}[z^{n-k}]f_r(z)^k,
                                                       \label{eq:S-coefficient}\\
 \Cnum rnk&=\frac{(n+(r-1)k)!}{k!}[z^{n-k}]g_r(z)^k.
                                                       \label{eq:C-coefficient}
\end{align}
No convergence assertion is needed, because each coefficient uses finitely
many terms; only degrees through $n-k$ need be retained.
These formulas are useful independent checks because they do not use the
two-term recurrence, and they reproduce the first-column and diagonal
formulas immediately. They test a different thing from the block-size
profile formula \eqref{eq:profile}: the generating function on one side,
the raw combinatorial definition on the other. \dup{D4}

\subsection{Row-polynomial operators}

Let $\vartheta=x\,d/dx$, and write
$T^{(r)}_n(x)=\sum_kT^{(r)}_{n,k}x^k$ and
$C^{(r)}_n(x)=\sum_kC^{(r)}_{n,k}x^k$.
Summing the recurrences, with multiplication by $x$ performed \emph{after}
the differential operator, gives
\begin{align}
 T^{(r)}_n(x)
 &=xP_d(n+d-1+d\vartheta)T^{(r)}_{n-1}(x)
                         +\vartheta T^{(r)}_{n-1}(x),\label{eq:Toperator}\\
 C^{(r)}_n(x)
 &=xP_d(n+d-1+d\vartheta)C^{(r)}_{n-1}(x)
                    +(n-1+d\vartheta)C^{(r)}_{n-1}(x).
                                                       \label{eq:Coperator}
\end{align}
Here a polynomial in $\vartheta$ acts diagonally on monomials:
$P_d(n+d-1+d\vartheta)x^j=P_d(n+d-1+dj)x^j$.
Thus the order and indexing of the operator are unambiguous.
These formulas offer a compact generating-polynomial description, but the
proof above does not require an analysis of polynomial zeros --- which, by
\cref{ex:nonreal}, would in any case be an analysis of nonreal zeros.

\subsection{Small rows}

Table~\ref{tab:rows} displays a few of the rows covered by the conjectural
cases. Empty places denote entries outside the triangle, not omitted data.
The zero column has been suppressed.

\begin{table}[htbp]
\centering
\small
\renewcommand{\arraystretch}{1.2}
\begin{tabular}{ccrrrr}
\toprule
$r$ & $n$ & $S_{n,1}^{(r)}$ & $S_{n,2}^{(r)}$ & $S_{n,3}^{(r)}$ & $S_{n,4}^{(r)}$\\
\midrule
3 & 2 & 1 & 10 & &\\
3 & 3 & 1 & 35 & 280 &\\
3 & 4 & 1 & 91 & 2100 & 15400\\
\addlinespace
4 & 2 & 1 & 35 & &\\
4 & 3 & 1 & 126 & 5775 &\\
4 & 4 & 1 & 336 & 45045 & 2627625\\
\addlinespace
5 & 2 & 1 & 126 & &\\
5 & 3 & 1 & 462 & 126126 &\\
5 & 4 & 1 & 1254 & 1009008 & 488864376\\
\bottomrule
\end{tabular}
\caption{Higher-order subset rows, computed exactly from \eqref{eq:Srec}
and checked independently using \eqref{eq:S-coefficient} and
\eqref{eq:profile}.}\label{tab:rows}
\end{table}

For $r=3$, the first six subset row sums, starting at $n=0$, are
\[
 1,\ 1,\ 11,\ 316,\ 17592,\ 1612206.
\]
For $r=4$ they are
\[
 1,\ 1,\ 36,\ 5902,\ 2673007,\ 2582136349,
\]
and for $r=5$,
\[
 1,\ 1,\ 127,\ 126589,\ 489874639,\ 5201522272948.
\]
The archive includes row sums through $n=30$ and complete small triangles
through $n=12$ for both kinds and $1\le r\le8$. These are reproducible
integer-sequence data; no assertion is made about whether each sequence
already has a separate OEIS entry.

\section{Exact computational audit}\label{sec:verification}

The mathematical proof is finite algebra plus induction over all rows.
The computations serve two distinct purposes: checking those finite
algebraic certificates, and testing for implementation or indexing mistakes
in the number arrays. Neither numerical root finding nor floating-point
inequalities are used anywhere. Two independent toolchains are supplied,
so that no single implementation --- and no single computer algebra
version --- is the sole witness of any identity. \dup{D5}

\subsection{Polynomial certificates without a computer algebra system}

The script \texttt{certificates.py} implements elementary polynomial
addition, multiplication, scalar division, and substitution over $\Q$.
Polynomials are tuples of rational coefficients in ascending degree.
All identities are checked by exact equality of the coefficient tuples,
not by evaluation at finitely many points. The recorded run verified
51 identities or shifted-coefficient positivity certificates, as well as
24 leading-coefficient checks for $d=1,\ldots,8$.

In particular it verifies \eqref{eq:MminusL}, \eqref{eq:main-factor},
the three coefficientwise identities that together constitute
\eqref{eq:SOS}, \eqref{eq:positive-cubic},
\eqref{eq:cycle4-discriminant}, \eqref{eq:cycle5-discriminant}, and
\eqref{eq:ratio-gap}; it also checks every polynomial listed in
\cref{sec:polynomial-table}. The shifted polynomials used for positivity
have nonnegative coefficients and positive constant term, so their
positivity on a nonnegative half-line is immediate. The full coefficient
arrays are written to \texttt{data/certificates.json}.

\subsection{The same certificates under SymPy}

The script \texttt{code/verify\_certificates.py} performs 24 named exact
checks with SymPy, expanding polynomial differences over the rationals and
requiring exact zero. Its list includes the universal four-term
decomposition \eqref{eq:slack-identity}, the lower-order expansions of
$A_d$ and $H_d$, the order-three and order-four boundary bounds, the
mixed-term remainder \eqref{eq:MminusL}, the discriminant factorization
\eqref{eq:main-factor}, the square certificate \eqref{eq:SOS} in one piece,
the positive quadratic \eqref{eq:Bpositive} and shifted cubic
\eqref{eq:positive-cubic}, the cycle mixed term and obstruction
\eqref{eq:cycle5-discriminant}, and the sharpness polynomial
\eqref{eq:ratio-gap}. Its report is \texttt{data/certificate\_checks.json},
which records the Python and SymPy versions actually used.
The two checkers are deliberately not a single program: the
CAS-free path certifies every identity in this manuscript on its own,
and the SymPy path checks the bivariate identity \eqref{eq:SOS} directly
rather than through its three coefficientwise components.

\subsection{Independent integer-array checks}

The script \texttt{verify.py} streams rows using \eqref{eq:Srec} and
\eqref{eq:Crec}. With its default settings, it checks all $n\le200$ for
$1\le r\le5$, both subset and cycle triangles. For each of the ten
triangles there are
\[
 \sum_{n=3}^{200}(n-2)=19701
\]
positive-support interior inequalities, giving $197010$ exact strict
inequalities in total. All passed. The order-five subset lower bound
\eqref{eq:quantitative} was also checked at all $19701$ such indices.
The order-five cycle run is explicitly marked \emph{experimental} in the
output; it is not part of the proved theorem.

The script \texttt{code/verify\_rows.py} performs a complementary scan
that is deeper but narrower: all $n\le300$ for the five subset triangles,
that is $44850$ internal inequalities for each order and $224250$ in
total. Neither scan range contains the other, and both are run. \dup{D6}
That script also evaluates the \emph{full} six-term certificate
\eqref{eq:full-certificate}, not merely its sign, at all $1711$ pairs with
$3\le n\le60$ and $2\le k\le n-1$, checking the identity and each
summand's sign separately; it checks the third-row closed form
\eqref{eq:row3} for $1\le r\le100$; and it locates the first failure for
each order $r=6,\ldots,20$, recording full triples and defects in
\texttt{data/row\_checks.json}. The counterexample table of
\texttt{data/verification.json} covers $r=6,\ldots,12$ and is retained
beside it rather than replaced. \dup{D10}

As structurally independent entry checks, neither of which uses the
recurrence, \texttt{verify.py} computes \eqref{eq:S-coefficient} and
\eqref{eq:C-coefficient} by truncated convolution of rational coefficient
lists, comparing every entry for $n\le12$, $1\le r\le8$ and both kinds
($1456$ entries, all agreeing), while \texttt{code/verify\_rows.py}
enumerates unordered block-size profiles and evaluates \eqref{eq:profile}
directly ($728$ entries for $1\le r\le8$, $0\le n\le12$, all agreeing).
First-column formulas, diagonal formulas, the exact counterexamples, and
the operator obstruction of \cref{sec:obstruction} are also checked. \dup{D4}

\begin{table}[htbp]
\centering
\small
\renewcommand{\arraystretch}{1.2}
\begin{tabular}{@{}p{0.56\textwidth}rr@{}}
\toprule
Check & Number & Result\\
\midrule
CAS-free polynomial identities and positivity certificates & 51 & Pass\\
CAS-free leading-coefficient checks, $d=1,\ldots,8$ & 24 & Pass\\
SymPy exact identities and sign certificates & 24 & Pass\\
Strict interior inequalities, ten triangles, $n\le200$ & 197010 & Pass\\
Quantitative fifth-order bounds, $n\le200$ & 19701 & Pass\\
Strict row inequalities, five subset triangles, $n\le300$ & 224250 & Pass\\
Independent generating-function entries, both kinds & 1456 & Pass\\
Independent block-profile entry counts & 728 & Pass\\
Full six-term certificate evaluations & 1711 & Pass\\
Third-row formula, orders $1$ through $100$ & 100 & Pass\\
First-failure tests, orders $6$ through $20$ & 15 & Pass\\
\bottomrule
\end{tabular}
\caption{The verification actually executed for this merged report.}
\label{tab:checks}
\end{table}

Deb and Sokal already report finite log-concavity testing through
$n=1000$ for the cases $r=3,4,5$ \cite[Appendix~C, Tables~1--2]{DebSokal}.
Neither of our smaller independent runs is therefore a claim to improve
that computational range. Their purpose is to audit these implementations
and the newly derived quantitative bound; the all-$n$ result is supplied by
the proof.

\subsection{Reproduction and trust boundary}

The standard-library scripts require only Python 3.9 or later
(\texttt{code/verify\_rows.py} uses Python 3.10 syntax) and no third-party
package:
\begin{lstlisting}
python certificates.py
python verify.py
python code/verify_rows.py --max-n 300 --output-dir data
\end{lstlisting}
The SymPy checker is the one component with an external dependency; the run
recorded here used Python 3.14.4 with SymPy 1.14.0, the version pinned in
\texttt{requirements.txt} (the second source draft recorded the same SymPy
version under Python 3.13.5):
\begin{lstlisting}
python -m pip install -r requirements.txt
python code/verify_certificates.py --output data/certificate_checks.json
\end{lstlisting}
A longer finite regression run is available explicitly:
\begin{lstlisting}
python verify.py --max-n 500
\end{lstlisting}
No network request, hidden data file, or proprietary computer algebra
service is required, and every identity in this manuscript is verified on
the CAS-free path alone. The assertions are implemented as explicit checks
and are not disabled by Python's \texttt{-O} flag.

The recurrence computation uses $O(N^2)$ integer-arithmetic operations and
stores $O(N)$ entries when rows are streamed, for fixed order $r$ and
maximum row $N$. These are arithmetic-operation counts, not bit-complexity
claims: the integers grow with $N$, and \texttt{data/row\_checks.json}
records the largest entry bit lengths actually exercised (for example
$9556$ bits at order five, $n=300$). The independent convolution and
block-profile checks are intentionally restricted to small rows, where
exact arithmetic is cheap and implementation independence is more valuable
than speed.

The scripts are not a proof-assistant formalization. Their trust boundary
is ordinary Python integer/rational arithmetic, the small supplied
polynomial implementation, and --- for the SymPy checker alone --- that
library's expansion routines. The paper's proofs do not logically depend on
executing the scripts: all identities, factor signs, support conditions,
and induction steps are written out.

\section{Conclusion}\label{sec:conclusion}

The chosen problem has a clean resolution: the higher-order Stirling
subset triangle has log-concave rows exactly for orders one through five.
The proof is short at its core. Normalize the recurrence, keep all terms
in the local quadratic bound, use the triangular constraint
$t-5k=n-k-1\ge0$, and apply one identity, \eqref{eq:SOS}, which converts a
negative mixed coefficient from an obstruction into a controlled term of a
positive quadratic. A single sixth-order counterexample plus a monotone
factorial ratio supplies the sharp converse.

The method also settles the third- and fourth-order cycle cases and yields
a continuous interpolating family, and the asymptotics of
\cref{sec:asymptotics} explain why the mechanism must change exactly
between orders five and six. Its failure at the fifth-order cycle
operator is genuine and explicitly certified, not merely a failure to
find a convenient factorization. Any extension to that remaining case
must exploit a stronger property of the actual preceding rows. The
present report separates that open continuation from the completed
subset-number proof. Independent mathematical review and a formal
proof-assistant development would be appropriate further checks.

\appendix
\section{Polynomial identities in one place}\label{sec:polynomial-table}

This appendix collects the low-degree expressions used in the proof.
The independent variable is $t$, and $H_d(t,k)$ is affine in $k$.
The notation matches the main text; the code uses \texttt{C0} and
\texttt{C1} for the constant and linear-in-$k$ parts of $H_d$.

\subsection{Pure-square coefficient}
\begin{align*}
 A_0(t)&=0,\\
 A_1(t)&=1,\\
 A_2(t)&=4(2t^2-2t-3),\\
 A_3(t)&=9(3t^4-12t^3-9t^2+42t+40),\\
 A_4(t)&=32(2t^6-18t^5+17t^4+168t^3-55t^2-834t-630).
\end{align*}

\subsection{Subset mixed coefficient and its lower bound}
\begin{align*}
 H_0(t,k)&=0,\\
 H_1(t,k)&=2,\\
 H_2(t,k)&=4(-2k+2t-1),\\
 H_3(t,k)&=6(-9kt+9k+3t^2-6t+11),\\
 H_4(t,k)&=16(-12kt^2+36kt-54k+2t^3-9t^2+43t-51).
\end{align*}
The substitution $k=t/(d+1)$ gives
\begin{align*}
 L_0(t)&=0,\qquad L_1(t)=2,\\
 L_2(t)&=\frac43(4t-3),\\
 L_3(t)&=\frac32(3t^2-15t+44),\\
 L_4(t)&=-\frac{16}{5}(2t^3+9t^2-161t+255)=-\frac{16}5\alpha(t).
\end{align*}

\subsection{Cycle mixed coefficient}
\begin{align*}
 M_0(t)&=0,\qquad M_1(t)=2,\\
 M_2(t)&=8(t-1),\\
 M_3(t)&=-18(3t-11),\\
 M_4(t)&=-32(t-2)(2t^2+4t-51).
\end{align*}

\subsection{The three fifth-order auxiliaries}
\begin{align*}
 \alpha(t)&=2t^3+9t^2-161t+255,\\
 \beta(t)&=2t^2-6t+9=2(t-3/2)^2+9/2,\\
 \gamma(t)&=7t^3-31t^2-126t+1080,\qquad
 \gamma(10+s)=7s^3+179s^2+1354s+3720.
\end{align*}
Also $\Delta_4(t)=96\beta(t)$, which is the slope identity behind
\eqref{eq:MminusL}.

\subsection{A direct expansion check for the main discriminant}

For readers checking \eqref{eq:main-factor} without software, both sides
expand to
\begin{align*}
 4A_4-L_4^2={}&\frac1{25}\bigl(
 5376t^6-66816t^5+198528t^4+1018368t^3\\
 &\hspace{20mm}{}-7986816t^2+18351360t-18662400\bigr).
\end{align*}
This displayed coefficient form is itself included in the certificate
checks. The factored form is preferable for proving its sign, because it
separates the domain condition $t\ge10$ from the algebraic identity.
\dup{D1}

\section{Deliberate redundancies}\label{sec:redundancy}

Merging the two source packages removed duplicated prose but deliberately
kept every place where the two drafts checked the same fact in genuinely
different ways. Collapsing any row below would lose an independent check.

\begin{center}
\small
\renewcommand{\arraystretch}{1.2}
\begin{longtable}{@{}>{\raggedright\arraybackslash}p{0.07\textwidth}
>{\raggedright\arraybackslash}p{0.38\textwidth}
>{\raggedright\arraybackslash}p{0.47\textwidth}@{}}
\toprule
Tag & Kept in more than one form & What the duplication buys\\
\midrule
\endhead
D1 & The fifth-order certificate as the single identity \eqref{eq:SOS}, as
the two-step discriminant plus completed square
\eqref{eq:main-factor}--\eqref{eq:hard-square}, and as the expanded
coefficient form in \cref{sec:polynomial-table}. &
The identity is the cleanest to audit, the two-step form is the one that
generalizes to the cycle case and to arbitrary $d$, and the expanded form
is the only one checkable entirely by hand.\\
D2 & Two routes to $H_4\ge L_4$: the exact identity \eqref{eq:MminusL},
and the affine-in-$k$ monotonicity of \eqref{eq:Hdecreasing} via
$\Delta_d>0$ (\cref{lem:falling}, \cref{rem:tworoutes}). &
The identity is the shorter proof at $d=4$; the convexity route is the
only one that survives to general $d$, and \cref{thm:interpolation} and
\cref{sec:asymptotics} both need it.\\
D3 & Two derivations of $S^{(6)}_{4,2}$ and $S^{(6)}_{4,3}$: block-size
pattern and recurrence (\cref{sec:sharpness}). &
This is the single load-bearing numeric fact of the converse; a
transcription error in either derivation would be caught by the other.\\
D4 & Two entry-verification routes that avoid the recurrence: truncated
generating-function convolution \eqref{eq:S-coefficient} and block-size
profiles \eqref{eq:profile}. &
One tests the generating function, the other the raw combinatorial
definition; an indexing error common to both is very unlikely.\\
D5 & Two verification toolchains: \texttt{certificates.py} (own exact
rational polynomial arithmetic, no CAS) and
\texttt{code/verify\_certificates.py} (SymPy). &
The shared identities are confirmed by two independent implementations,
and the package is not hostage to one CAS version.\\
D6 & Two finite scan ranges, neither containing the other: $n\le200$
across ten subset and cycle triangles, and $n\le300$ across the five
subset triangles. &
The first covers the cycle arrays and the quantitative bound; the second
goes deeper in $n$ and exercises much larger integers.\\
D7 & Two witnesses for the fifth-order cycle obstruction: the
factorization \eqref{eq:cycle5-discriminant} and the exact numerical input
at $n=5$, $q=7350$. &
The algebraic witness shows the criterion fails identically for $t\ge10$;
the numerical one exhibits an actual log-concave input whose image is not
log-concave.\\
D8 & Two attributions of the coefficient-sign condition: Sagan
\cite[Theorem~1]{Sagan} and Shankar \cite[Theorem~1.2]{Shankar}. &
The same sufficient condition is recorded in two independent sources, and
each framing of how the present criterion relaxes it is worth keeping.\\
D9 & Two presentations of the sharpness ratio: $F(r)$ decreasing and
$R(r)=1/F(r)$ increasing, with both ratio expressions in
\eqref{eq:F-ratio}. &
Each ratio expression is a transcription check on the other, and both
routes reach the same difference polynomial \eqref{eq:ratio-gap}.\\
D10 & Two recorded counterexample tables: \texttt{data/verification.json}
for $r=6,\ldots,12$ and \texttt{data/row\_checks.json} for
$r=6,\ldots,20$. &
The union is kept; the second extends the range and adds bit-length
instrumentation, the first adds the cycle-operator obstruction record.\\
D11 & Two shape consequences: ratio monotonicity with the $q$-weighted
log-concave mass function, and the balanced-product inequalities
\eqref{eq:balanced}. &
Neither statement implies the other as stated: the balanced products are
stronger as inequalities, the weighted-distribution remark is about a
different object.\\
D12 & Two scope fences: the result-to-source map in \cref{sec:status} and
the nonreal-rootedness \cref{ex:nonreal}. &
The first guards against overclaiming which conjecture clauses are
settled; the second rules out the real-rootedness and total-positivity
readings of the result.\\
\bottomrule
\end{longtable}
\end{center}

\section{A minimal symbolic check}\label{sec:minimal-check}

The following compact calculation verifies the central equality
\eqref{eq:SOS} alone, in the letters of this manuscript. The two full
checkers in the archive verify all the auxiliary identities as well, and
\texttt{certificates.py} does so without SymPy.
\begin{lstlisting}[language=Python]
import sympy as s

t, k, u, v = s.symbols("t k u v")
P4 = t*(t-1)*(t-2)*(t-3)
pm, pp = P4.subs(t, t-4), P4.subs(t, t+4)
A4 = P4*P4 - pm*pp
H4 = 2*k*P4 - (k+1)*pm - (k-1)*pp
alpha = 2*t**3 + 9*t*t - 161*t + 255
beta = 2*t*t - 6*t + 9
gamma = 7*t**3 - 31*t*t - 126*t + 1080
lhs = 25*(A4*u*u + H4*u*v + v*v)
rhs = ((5*v - 8*alpha*u)**2 + 96*(t-5)*beta*gamma*u*u
       + 480*beta*(t - 5*k)*u*v)
assert s.expand(lhs - rhs) == 0
\end{lstlisting}
An exact polynomial identity is only one part of the proof. Its
application also needs the combinatorial recurrence, the support and
boundary analysis, the preceding row's log-concavity, and the domain
$t\ge5k\ge10$. Those mathematical hypotheses are proved explicitly in
the main text rather than delegated to this code.

\section{Source and status notes}\label{sec:status}

The problem identification uses the specific version
\texttt{arXiv:2507.18959v1}. The conjecture and recurrence were checked in
both its HTML and PDF renderings. On the report date the arXiv abstract
page showed only the July 2025 version in its submission history.
Searches for the exact paper identifier, the phrase ``higher-order Stirling''
with log-concavity, and relevant arXiv author/topic combinations did not
locate a subsequent resolution. The BCC30 statement and the recent
triangular-recurrence paper \cite{Shankar} were checked separately.
Both source drafts performed such a check independently, on 19 and on
20 September 2026; neither located a later resolution. This is a documented
bounded search, not proof of absence from all literature, and it does not
establish priority.

The underlying combinatorial objects, their recurrences, and the target
conjecture are from the cited literature. What the present report offers as
its own contribution is the fifth-order certificate and its two forms, the
induction, the quantitative bound, the sharpness argument, the cycle
extensions and obstruction, the interpolating family, the asymptotic
explanation, and the accompanying checks. No novelty is claimed for the
general two-variable preservation lemma, for the elementary identities, or
for the small counterexamples.

The mathematical claims established here can therefore be assessed
independently of any priority judgment. The result-to-source map is:
\begin{center}
\small
\renewcommand{\arraystretch}{1.25}
\begin{tabular}{>{\raggedright\arraybackslash}p{0.34\textwidth}>{\raggedright\arraybackslash}p{0.56\textwidth}}
\toprule
Source assertion & Result in this manuscript\\
\midrule
Deb--Sokal 1.4(c), row log-concavity & All stated orders $1\le r\le5$;
strictness and sharp converse also proved.\\
Deb--Sokal 1.3(c), row log-concavity & Orders $1\le r\le4$ proved;
order five not settled.\\
BCC30 Problem 30.9 & Cycle orders three and four proved;
order five not settled.\\
Other clauses of those conjectures & Not claimed as consequences of the results here.\\
\bottomrule
\end{tabular}
\end{center}

\begin{thebibliography}{9}
\bibitem{DebSokal}
Bishal Deb and Alan D. Sokal,
\emph{Higher-order Stirling cycle and subset triangles: Total positivity,
continued fractions and real-rootedness},
\href{https://arxiv.org/abs/2507.18959}{arXiv:2507.18959v1}, 2025.
See Lemma~1.1, Conjectures~1.3(c) and 1.4(c), and Appendix~C.

\bibitem{Cameron}
Peter Cameron (editor), \emph{Problems from BCC30},
\href{https://arxiv.org/abs/2409.07216}{arXiv:2409.07216v1}, 2024.
Problem~30.9, presented by Bishal Deb, concerns the cycle triangle and is
not a separate source for the subset conjecture.

\bibitem{Sagan}
Bruce E. Sagan, \emph{Inductive and injective proofs of log concavity results},
Discrete Mathematics \textbf{68} (1988), 281--292.
\href{https://doi.org/10.1016/0012-365X(88)90120-3}{doi:10.1016/0012-365X(88)90120-3}.
See Theorem~1 and its inductive proof.

\bibitem{Shankar}
Umesh Shankar, \emph{Log-concavity of rows of triangular arrays satisfying
a certain super-recurrence},
\href{https://arxiv.org/abs/2508.12467}{arXiv:2508.12467v1}, 2025.
See Theorem~1.2 for the coefficient-sign criterion, Theorems~1.3--1.4 for
its relatives, and Section~1 for the different $r$-Stirling family.

\bibitem{OEIS}
The OEIS Foundation Inc., \emph{The On-Line Encyclopedia of Integer Sequences},
\href{https://oeis.org/A134991}{entry A134991: Triangle of Ward numbers},
accessed 19 September 2026.
\end{thebibliography}
\end{document}
